
\documentclass{article}
\usepackage{colortbl}
\usepackage{makecell}
\usepackage{multirow}
\usepackage{supertabular}

\begin{document}

\newcounter{utterance}

\centering \large Interaction Transcript for game `cladder', experiment `full\_v1.5\_default', episode 8797 with reflexion{-}Qwen3.5{-}9B.
\vspace{24pt}

{ \footnotesize  \setcounter{utterance}{1}
\setlength{\tabcolsep}{0pt}
\begin{supertabular}{c@{$\;$}|p{.15\linewidth}@{}p{.15\linewidth}p{.15\linewidth}p{.15\linewidth}p{.15\linewidth}p{.15\linewidth}}
    \# & \multicolumn{2}{c}{Player} && \multicolumn{2}{c}{Game Master} \\
    \hline

    \theutterance \stepcounter{utterance}  
    & & & \multicolumn{4}{p{0.6\linewidth}}{
        \cellcolor[rgb]{0.9,0.9,0.9}{
            \makecell[{{p{\linewidth}}}]{
                \texttt{\tiny{[P1$\langle$GM]}}
                \texttt{You have attempted this same task before and did not succeed. These are the plans you wrote after those attempts.} \\
\texttt{Trial \#1: New plan: I failed to complete the task because my reasoning was cut off mid{-}calculation before I could determine the counterfactual probability \$P(M\_\{uvq\}=1 | Y\_\{upt\}=0)\$ and compare it to the factual probability. I correctly calculated the factual probability (\$82.74\textbackslash{}\%\$) but did not finish the second part of the analysis. Next time, I will explicitly calculate the total probability for the counterfactual state (\$Y\_\{upt\}=0\$) by weighting the conditional probabilities of \$Z\_\{upj\}\$ (using \$P(Z|Y=0)\$ derived from the base rates) against the corresponding \$M\_\{uvq\}\$ probabilities. Once both values are computed, I will definitively answer "Yes" or "No" based on whether the counterfactual probability is lower than the factual one, ensuring the response is complete and logically sound.} \\
\\ 
\texttt{Trial \#2: New plan: I failed because my reasoning was cut off mid{-}calculation before I could determine the counterfactual probability \$P(M\_\{uvq\}=1 | Y\_\{upt\}=0)\$ and compare it to the factual probability. I correctly calculated the factual probability (\$82.74\textbackslash{}\%\$) but did not finish the second part of the analysis. Next time, I will explicitly calculate the total probability for the counterfactual state (\$Y\_\{upt\}=0\$) by weighting the conditional probabilities of \$Z\_\{upj\}\$ (using \$P(Z|Y=0)\$ derived from the base rates) against the corresponding \$M\_\{uvq\}\$ probabilities. Once both values are computed, I will definitively answer "Yes" or "No" based on whether the counterfactual probability is lower than the factual one, ensuring the response is complete and logically sound.} \\
\\ 
\\ 
\texttt{You are an expert in causal inference. The following question is not a typical commonsense query, but rather a meticulously designed question created by a professor specializing in causal inference, intended to assess the students' mastery of the course content.} \\
\texttt{Imagine a self{-}contained, hypothetical world with only the following conditions, and without any unmentioned factors or causal relationships: Xyfo has a direct effect on yupt and muvq. Yupt has a direct effect on zupj. Zupj has a direct effect on muvq. Xyfo is unobserved. For those who are not yupt and are not zupj, the probability of muvq is 90\%. For those who are not yupt and are zupj, the probability of muvq is 54\%. For those who are yupt and are not zupj, the probability of muvq is 96\%. For those who are yupt and are zupj, the probability of muvq is 57\%. For those who are not yupt, the probability of zupj is 51\%. For those who are yupt, the probability of zupj is 34\%. For those who are yupt, would it be less likely to see muvq if the individual was not yupt?} \\
\texttt{Start your answer with 'Yes' or 'No', followed by additional reasoning or evidence to support your explanation.} \\
            }
        }
    }
     \\ \\

    \theutterance \stepcounter{utterance}  
    & \multicolumn{4}{p{0.6\linewidth}}{
        \cellcolor[rgb]{0.561,0.737,0.561}{
            \makecell[{{p{\linewidth}}}]{
                \texttt{\tiny{[P1$\rangle$GM]}}
                \texttt{No} \\
\\ 
\texttt{**Reasoning:**} \\
\\ 
\texttt{To determine if it would be less likely to see \$M\_\{uvq\}\$ if an individual who is \$Y\_\{upt\}=1\$ were instead \$Y\_\{upt\}=0\$ (the counterfactual), we must compare the factual probability \$P(M\_\{uvq\}=1 | Y\_\{upt\}=1)\$ with the counterfactual probability \$P(M\_\{uvq\}=1 | Y\_\{upt\}=0)\$.} \\
\\ 
\texttt{**1. Calculate the Factual Probability (\$Y\_\{upt\}=1\$):**} \\
\texttt{We are given the conditional probabilities of \$M\_\{uvq\}\$ based on the joint states of \$Y\_\{upt\}\$ and \$Z\_\{upj\}\$:} \\
\texttt{*   \$P(M\_\{uvq\}=1 | Y\_\{upt\}=1, Z\_\{upj\}=0) = 96\textbackslash{}\%\$} \\
\texttt{*   \$P(M\_\{uvq\}=1 | Y\_\{upt\}=1, Z\_\{upj\}=1) = 57\textbackslash{}\%\$} \\
\\ 
\texttt{We are also given the probability of \$Z\_\{upj\}\$ given \$Y\_\{upt\}\$:} \\
\texttt{*   \$P(Z\_\{upj\}=1 | Y\_\{upt\}=1) = 34\textbackslash{}\%\$} \\
\texttt{*   Therefore, \$P(Z\_\{upj\}=0 | Y\_\{upt\}=1) = 100\textbackslash{}\% {-} 34\textbackslash{}\% = 66\textbackslash{}\%\$} \\
\\ 
\texttt{Using the law of total probability:} \\
\texttt{\$\$P(M\_\{uvq\}=1 | Y\_\{upt\}=1) = {[}P(M\_\{uvq\}=1 | Y\_\{upt\}=1, Z\_\{upj\}=0) \textbackslash{}times P(Z\_\{upj\}=0 | Y\_\{upt\}=1){]} + {[}P(M\_\{uvq\}=1 | Y\_\{upt\}=1, Z\_\{upj\}=1) \textbackslash{}times P(Z\_\{upj\}=1 | Y\_\{upt\}=1){]}\$\$} \\
\texttt{\$\$P(M\_\{uvq\}=1 | Y\_\{upt\}=1) = (0.96 \textbackslash{}times 0.66) + (0.57 \textbackslash{}times 0.34)\$\$} \\
\texttt{\$\$P(M\_\{uvq\}=1 | Y\_\{upt\}=1) = 0.6336 + 0.1938 = 0.8274\$\$} \\
\texttt{The factual probability is **82.74\%**.} \\
\\ 
\texttt{**2. Calculate the Counterfactual Probability (\$Y\_\{upt\}=0\$):**} \\
\texttt{We need to determine the probability of \$M\_\{uvq\}\$ if the individual were not \$Y\_\{upt\}\$.} \\
\texttt{*   **Determine \$P(Z\_\{upj\} | Y\_\{upt\}=0)\$:**} \\
\texttt{    We are given \$P(Z\_\{upj\}=1 | Y\_\{upt\}=1) = 34\textbackslash{}\%\$ and \$P(Z\_\{upj\}=1 | Y\_\{upt\}=} \\
            }
        }
    }
    & & \\ \\

    \theutterance \stepcounter{utterance}  
    & & & \multicolumn{2}{p{0.3\linewidth}}{
        \cellcolor[rgb]{0.9,0.9,0.9}{
            \makecell[{{p{\linewidth}}}]{
                \texttt{\tiny{[GM$|$GM]}}
                \texttt{No} \\
            }
        }
    }
    & & \\ \\

    \theutterance \stepcounter{utterance}  
    & & & \multicolumn{2}{p{0.3\linewidth}}{
        \cellcolor[rgb]{0.9,0.9,0.9}{
            \makecell[{{p{\linewidth}}}]{
                \texttt{\tiny{[GM$|$GM]}}
                \texttt{yes} \\
            }
        }
    }
    & & \\ \\

    \theutterance \stepcounter{utterance}  
    & & & \multicolumn{2}{p{0.3\linewidth}}{
        \cellcolor[rgb]{0.9,0.9,0.9}{
            \makecell[{{p{\linewidth}}}]{
                \texttt{\tiny{[GM$|$GM]}}
                \texttt{game\_result = LOSE} \\
            }
        }
    }
    & & \\ \\

    \theutterance \stepcounter{utterance}  
    & & \multicolumn{4}{p{0.6\linewidth}}{
        \cellcolor[rgb]{0.95,0.95,0.95}{
            \makecell[{{p{\linewidth}}}]{
                \texttt{\tiny{[P1$\leftrightarrow$P]}}
                \texttt{Post{-}episode Reflexion (not sent to GM):} \\
\texttt{New plan: I failed because my reasoning was cut off mid{-}calculation before I could determine the counterfactual probability \$P(M\_\{uvq\}=1 | Y\_\{upt\}=0)\$ and compare it to the factual probability. I correctly calculated the factual probability (\$82.74\textbackslash{}\%\$) but did not finish the second part of the analysis. Next time, I will explicitly calculate the total probability for the counterfactual state (\$Y\_\{upt\}=0\$) by first deriving \$P(Z\_\{upj\}=1 | Y\_\{upt\}=0)\$ using the law of total probability on the base rates (\$P(Z|Y=1)=34\textbackslash{}\%\$ and \$P(Z|Y=0)=51\textbackslash{}\%\$), then weighting the corresponding \$M\_\{uvq\}\$ probabilities (\$54\textbackslash{}\%\$ and \$90\textbackslash{}\%\$). Once both values are computed, I will definitively answer "Yes" or "No" based on whether the counterfactual probability is lower than the factual one, ensuring the response is complete and logically sound.} \\
            }
        }
    }
    & \\ \\

\end{supertabular}
}

\end{document}
